\section*{Chapter \ref{ch:ml:data-and-feature-stores} --- Data and Feature Stores}

\begin{solution}{exo:ml:data-and-feature-stores:1}
The \omterm{def:ml:data-and-feature-stores:store}{online store}: trades stamped in $(09{:}59{:}37.300, 10{:}00{:}07.300]$. The bar pipeline: trades in the thirty one-second
bars ending at 10:00:08, that is $(09{:}59{:}38, 10{:}00{:}08]$: it misses 0.7~s of the past and includes 0.7~s of the future.
\end{solution}

\begin{solution}{exo:ml:data-and-feature-stores:2}
$1 - 0.999^n\ge0.95$ gives $n\ge\ln0.05/\ln0.999 = 2\,994$ decisions: about 3\,000, far more than a five-decision check, which is
why parity checks should also run on full days in batch.
\end{solution}

\begin{solution}{exo:ml:data-and-feature-stores:3}
The volume and count features are sums of integers, represented exactly in float32 up to $2^{24}$; the imbalances and the
weighted mid are ratios with 24 bits of precision in float32 against 53 in float64, so they change in the seventh or eighth
significant digit.
\end{solution}

\begin{solution}{exo:ml:data-and-feature-stores:4}
A bar's end is up to a second after the decision: for a one-second forecast the leaked interval is a large part of the target
window, and the research model learned to use it; for a five-second forecast it is at most a fifth of the window and the
information is diluted.
\end{solution}

\begin{solution}{exo:ml:data-and-feature-stores:5}
The message after a trade is usually a limit order added or cancelled deep in the book: its effect on the target is small and
it adds noise to features the model relies on, so the research model is slightly worse than one trained on clean features. If
the next message were typically the mid-changing one (a thin book, a fast market), the look-ahead would inflate the backtest
instead; the sign of a leak's effect is not knowable in advance.
\end{solution}

\begin{solution}{exo:ml:data-and-feature-stores:6}
Equal moments do not mean equal values: one-second bars shift values in time without changing their distribution, and a skew
that affects some decisions can leave the daily averages intact. Compare the values themselves, decision by decision.
\end{solution}

\begin{solution}{exo:ml:data-and-feature-stores:7}
Trained on the production engine's replay, the model's live IC is 0.411, the clean result, against 0.360 for the model trained on
bar features: everything the skew cost is recovered. It costs a replay of the full message history through the online engine to
build the \omterm{def:ml:why-financial-machine-learning-is-different:sets}{training set}, slower than the bar pipeline, which is exactly what a \omterm{def:ml:data-and-feature-stores:store}{feature store}'s shared definition makes routine.
\end{solution}

\begin{solution}{exo:ml:data-and-feature-stores:8}
The error of a float64 sum of $n$ terms in any order is bounded by about $n\,u\sum|x_i|$ with $u = 2^{-53}$; use a tolerance
proportional to $\sum|x_i|$ times a small multiple of $n\,u$, or make both computations exact (integer inputs, compensated
summation, Book~4, chapter~25) and demand equality.
\end{solution}

\begin{solution}{pb:ml:data-and-feature-stores:1}
\textbf{Part I.}
\begin{enumerate}
\item Name, input, aggregation, window, clock and version; as data so that both stores, the documentation and the tests read the
same declaration.
\item Offline: vectorised over a history from cumulative sums and as-of lookups; online: running windows updated message by
message. They agree exactly at all 7\,714 decisions.
\item The time of the last message it read; for a fundamental, its delivery time to the firm.
\item The age of a value's newest input; trade prints arriving 200~ms late made the trade features stale.
\end{enumerate}
\textbf{Part II.}
\begin{enumerate}[resume]
\item One event ahead (all book features, slightly the windows), one-second bars (every feature, most strongly the book's last
values), float32 storage (the ratio features, invisibly), late trade prints (the three trade features).
\item The first two are research bugs; the last two are production conditions.
\item 100\%, 96\%, 92\% and 100\% of decisions (at tolerance $10^{-9}$).
\item At $10^{-9}$ it flags the harmless float32 storage; at $10^{-6}$ it ignores it and still flags the others.
\end{enumerate}
\textbf{Part III.}
\begin{enumerate}[resume]
\item Backtest and live: clean 0.411 and 0.411; one event ahead 0.395 and 0.418; one-second bars 0.438 and 0.360; float32 0.411
and 0.411; late trades 0.411 and 0.411.
\item Its features saw part of the target window; the model weighted them; live they carry no such information.
\item The same skews move the IC by less than 0.01.
\item A live IC well below the backtest's from the first day, stable thereafter: the signature of a pipeline difference, not of
decay.
\end{enumerate}
\textbf{Part IV.}
\begin{enumerate}[resume]
\item \emph{Same features twice.} Research on one-second bars inflates the one-second forecast's backtest IC to 0.438 and
delivers 0.360 live against 0.411 for the clean pipeline, with P\&L per decision falling from 0.089 to 0.078 ticks; a
five-decision parity check catches each of the four planted skews every time.
\item The bar misalignment: it inflates research and hurts production, and nothing in either alone reveals it.
\item As a new version computed alongside the old, backfilled offline, compared online, and switched only when models are
retrained on it.
\item Stop trading the affected models or fall back, find the feature and the cause, fix, recompute, and retrain if the
training data were wrong.
\item Through the \omterm{def:ml:alternative-data-pipelines:ingest}{ingestion pipeline} of chapter~15, with first-seen times as knowledge times, then as ordinary definitions.
\item Definitions, versions, parity history and freshness records for every feature the model uses.
\item Engineering time and discipline; worth it as soon as two models share features or one model trades live.
\item That the model sees in production what it saw in training, as of the moment it decides.
\end{enumerate}
\end{solution}

\begin{solution}{iq:ml:data-and-feature-stores:1}
A difference between the features a model trained on and those it gets live: different code paths, different clocks or bars,
late or missing inputs, different arithmetic, look-ahead in the research join.
\iqlookfor{the definition and three distinct causes.}
\end{solution}

\begin{solution}{iq:ml:data-and-feature-stores:2}
Declarative definitions; an online engine updating running windows per message with bounded latency; an \omterm{def:ml:data-and-feature-stores:store}{offline store}
replaying the same engine or an equivalent vectorised computation over history, stamped with knowledge times; point-in-time
retrieval for training; daily parity and freshness monitoring; versioning.
\iqlookfor{shared definitions, point-in-time offline, parity and freshness.}
\end{solution}

\begin{solution}{iq:ml:data-and-feature-stores:3}
Replay the live period's messages through the research pipeline and compare features with what production logged; retrain on
production's replayed features and see whether the backtest falls to the live IC; check freshness logs.
\iqlookfor{a replay comparison before any modelling explanation.}
\end{solution}

\begin{solution}{iq:ml:data-and-feature-stores:4}
For each (entity, decision time), take each feature's latest value with knowledge time at or before the decision; implemented
by as-of joins on sorted times, never by joining on calendar keys.
\iqlookfor{knowledge time and the as-of rule.}
\end{solution}

\begin{solution}{iq:ml:data-and-feature-stores:5}
Record each input's last arrival and each feature's newest input; alert when their age exceeds a bound set by the feature's
window and the model's horizon; have the model refuse or fall back on stale features rather than trade on them.
\iqlookfor{measured age, bounds and a defined behaviour when stale.}
\end{solution}

\begin{solution}{iq:ml:data-and-feature-stores:6}
Make the inputs integers (ticks, shares) or fixed point, so sums are exact; or use the same summation order and compensated
summation in both; and test on a replay that every value matches exactly.
\iqlookfor{exact arithmetic or controlled order, and a replay test.}
\end{solution}
